\documentclass[12pt]{article}
\def\activityid{F07-M-v3}\usepackage[letterpaper,margin=.65in,top=.60in,bottom=.65in]{geometry}
\usepackage[T1]{fontenc}\usepackage[default]{sourcesanspro}
\usepackage{microtype,tikz,array,tabularx,fancyhdr,amsmath}\usepackage[hidelinks]{hyperref}
\usetikzlibrary{arrows.meta,calc}
\setlength{\parindent}{0pt}\setlength{\parskip}{7pt}\setlength{\headheight}{0pt}\setlength{\footskip}{26pt}\setlength{\emergencystretch}{2em}
\pagestyle{fancy}\fancyhf{}\renewcommand{\headrulewidth}{0pt}\renewcommand{\footrulewidth}{.4pt}
\fancyfoot[L]{\scriptsize Bellingham Math Circle / Week 7 / \activityid}\fancyfoot[R]{\scriptsize \thepage}
\definecolor{ink}{gray}{.12}\definecolor{soft}{gray}{.95}\color{ink}
\newcommand{\PageTitle}[2]{{\fontsize{10}{12}\selectfont\bfseries WEEK 7\quad GAME LAB\hfill #1\par}\vspace{3pt}{\fontsize{26}{29}\selectfont\bfseries #2\par}\vspace{2pt}}
\newcommand{\Subhead}[1]{\par\vspace{3pt}{\large\bfseries #1}\par}
\newcommand{\RuleBox}[1]{\par\begingroup\setlength{\fboxsep}{8pt}\colorbox{soft}{\parbox{\dimexpr\linewidth-16pt\relax}{#1}}\endgroup\par}
\newcommand{\WriteBox}[2]{\par\textbf{#1}\par\vspace{2pt}\begin{tikzpicture}\draw[black!40,line width=.5pt] (0,0) rectangle (\dimexpr\linewidth-.5pt\relax,#2);\end{tikzpicture}\par}
\newcommand{\Code}[1]{\texttt{#1}}

\newcommand{\StudentHeader}[1]{{\fontsize{10}{12}\selectfont\bfseries Week 7 / Strategy lab / #1\par}\vspace{10pt}}
\newcommand{\Problem}[2]{\par\vspace{4pt}\textbf{Problem #1:} #2\par}
\newcommand{\RecordBox}[2]{\par #1\par\vspace{2pt}\begin{tikzpicture}\draw[black!40,line width=.5pt](0,0)rectangle(\dimexpr\linewidth-.5pt\relax,#2);\end{tikzpicture}\par}

\hypersetup{pdftitle={Week 7: Grades 2--3}}
\begin{document}
\StudentHeader{Grades 2--3}
\Problem{1}{Two players alternate taking 1, 3, or 4 counters. Taking 2 is not allowed. Never take more than remain. Taking the last counter wins. Play from 5, 6, and 7. Swap who starts and replay each pile. Record the moves in order; a partner may write while you play.}
\begin{center}\renewcommand{\arraystretch}{2.1}\begin{tabular}{c|c|p{2.2in}|p{1.4in}}Starting pile & First move & Next moves, in order & Winner: first or second\\\hline
5 &  &  & \\
5 &  &  & \\
6 &  &  & \\
6 &  &  & \\
7 &  &  & \\
7 &  &  & \\\end{tabular}\end{center}
\Problem{2}{Build a pile of 2, then a pile of 4. From each pile, try every allowed first move. Reset the pile before the next attempt. Record how many counters each move leaves.}
\begin{center}\renewcommand{\arraystretch}{2.5}\begin{tabular}{c|p{5.5in}}Pile & Allowed move / counters left\\\hline
2 & \\
4 & \\\end{tabular}\end{center}

\newpage
\StudentHeader{Grades 2--3}
\Problem{3}{Use moves 1, 3, and 4. A position is W if the player about to move has a winning strategy; it is L if every move lets the other player force a win. Put L at 0 because there is no move. Work from small piles upward. List all legal next piles. Mark W if at least one is L; mark L if all are W.}
\begin{center}\renewcommand{\arraystretch}{2.0}\begin{tabular}{c|p{2.1in}|c|p{1.8in}}Pile & All legal next piles & W or L & One move to L, if any\\\hline
0 & none & L & none\\
1 &  &  & \\
2 &  &  & \\
3 &  &  & \\
4 &  &  & \\
5 &  &  & \\
6 &  &  & \\
7 &  &  & \\\end{tabular}\end{center}
\Problem{4}{Build a nonempty pile you marked L, such as 2 or 7. Let your partner choose a legal move, then use your chart to reply. Reset and try every different first move. Record your replies.}
\RecordBox{Starting pile, each first move, and my replies:}{1.1in}

\newpage
\StudentHeader{Grades 2--3}
\Problem{5}{Use moves 1, 3, and 4. From 8 counters, try taking 1, 3, and 4. Reset each time. Use your chart to record what each move leaves and whether that next pile is W or L. Circle every winning first move.}
\begin{center}\renewcommand{\arraystretch}{1.9}\begin{tabular}{c|c|c}Take & Pile left & Its W/L label\\\hline
1 &  & \\
3 &  & \\
4 &  & \\\end{tabular}\end{center}
\Problem{6}{Do the same from 10. First decide the missing labels for 8 and 9, using all their legal moves. Record all winning moves from 10, not just the first one you find.}
\begin{center}\renewcommand{\arraystretch}{1.9}\begin{tabular}{c|c|c}Take & Pile left & Its W/L label\\\hline
1 &  & \\
3 &  & \\
4 &  & \\\end{tabular}\end{center}
\Problem{7}{Choose a starting pile from 8 through 14. Make a strategy card: record each possible first move and the pile it leaves. Circle moves that hand your opponent an L. Test the card in a game.}
\RecordBox{My pile and strategy card:}{1.3in}

\end{document}
